% ============================================================================
%  Reporte de Residencia Profesional
%  Proyecto: Metaverso Escolar TecNM — Backend, Panel Docente e Integración LTI
%  Autor: Daniel Octavio Ramírez Neri
%  Compilar con: tectonic residencia.tex   (motor XeTeX)
% ============================================================================
\documentclass[12pt,letterpaper,oneside]{report}

% ---- Datos editables de la portada (rellena los que falten) ----
\newcommand{\Institucion}{Tecnológico Nacional de México}
\newcommand{\Instituto}{Instituto Tecnológico de \rule{4.2cm}{0.4pt}}
\newcommand{\Carrera}{Ingeniería en Sistemas Computacionales}
\newcommand{\TituloProyecto}{Metaverso Escolar TecNM: backend de prácticas, panel docente e integración con Moodle vía LTI 1.3}
\newcommand{\Alumno}{Daniel Octavio Ramírez Neri}
\newcommand{\NumControl}{\rule{3.5cm}{0.4pt}}
\newcommand{\AsesorInterno}{\rule{5cm}{0.4pt}}
\newcommand{\AsesorExterno}{\rule{5cm}{0.4pt}}
\newcommand{\Periodo}{Enero -- Junio 2026}
\newcommand{\Lugar}{México}

\usepackage[spanish,es-tabla]{babel}
\usepackage{geometry}
\geometry{letterpaper,top=2.5cm,bottom=2.5cm,left=3cm,right=2.5cm}
\usepackage{fontspec}
\usepackage{graphicx}
\usepackage{xcolor}
\usepackage{booktabs}
\usepackage{longtable}
\usepackage{array}
\usepackage{enumitem}
\usepackage{listings}
\usepackage{fancyhdr}
\usepackage{titlesec}
\usepackage{setspace}
\usepackage{tikz}
\usetikzlibrary{shapes.geometric,arrows.meta,positioning}
\usepackage[hidelinks]{hyperref}
\hypersetup{
  colorlinks=true, linkcolor=black, citecolor=black,
  urlcolor=blue!60!black, pdftitle={Reporte de Residencia Profesional},
  pdfauthor={Daniel Octavio Ramirez Neri}
}

\definecolor{azultec}{RGB}{31,59,110}
\definecolor{griscodigo}{RGB}{245,246,250}
\definecolor{verdecodigo}{RGB}{0,128,96}

\onehalfspacing

% ---- Encabezado / pie ----
\pagestyle{fancy}
\fancyhf{}
\fancyhead[L]{\small\itshape Metaverso Escolar TecNM}
\fancyhead[R]{\small\thepage}
\renewcommand{\headrulewidth}{0.4pt}

% ---- Estilo de capítulos ----
\titleformat{\chapter}[hang]
  {\normalfont\Large\bfseries\color{azultec}}
  {\thechapter.}{0.6em}{}
\titlespacing*{\chapter}{0pt}{6pt}{16pt}

% ---- Estilo de código ----
\lstdefinestyle{codigo}{
  backgroundcolor=\color{griscodigo},
  basicstyle=\ttfamily\footnotesize,
  breaklines=true,
  frame=single,
  rulecolor=\color{gray!40},
  numbers=left,
  numberstyle=\tiny\color{gray},
  keywordstyle=\color{azultec}\bfseries,
  commentstyle=\itshape\color{gray},
  stringstyle=\color{verdecodigo},
  showstringspaces=false,
  columns=fullflexible,
  tabsize=2,
  captionpos=b,
}
\lstset{style=codigo}

\begin{document}

% =========================== PORTADA ===========================
\begin{titlepage}
\centering
\vspace*{0.5cm}
{\scshape\large \Institucion\par}
\vspace{0.2cm}
{\scshape \Instituto\par}
\vspace{1.2cm}

\rule{\linewidth}{0.5pt}\par
\vspace{0.4cm}
{\Huge\bfseries\color{azultec} Reporte de Residencia Profesional\par}
\vspace{0.4cm}
\rule{\linewidth}{0.5pt}\par
\vspace{1.4cm}

{\large\bfseries \TituloProyecto\par}
\vspace{2.2cm}

{\large Presenta:\par}
\vspace{0.2cm}
{\Large\bfseries \Alumno\par}
\vspace{0.2cm}
{\normalsize Número de control: \NumControl\par}
\vspace{1.6cm}

\begin{tabular}{ll}
Carrera: & \Carrera \\[4pt]
Asesor interno: & \AsesorInterno \\[4pt]
Asesor externo: & \AsesorExterno \\[4pt]
Periodo: & \Periodo \\
\end{tabular}

\vfill
{\large \Lugar, \Periodo\par}
\end{titlepage}

% =========================== PRELIMINARES ===========================
\pagenumbering{roman}

\chapter*{Agradecimientos}
\addcontentsline{toc}{chapter}{Agradecimientos}
A mi institución, a mis asesores y a mi familia, por el apoyo brindado durante el
desarrollo de esta residencia profesional. Este proyecto representa la
integración de los conocimientos adquiridos a lo largo de la carrera en un
sistema real, útil y extensible para la comunidad educativa.

\chapter*{Resumen}
\addcontentsline{toc}{chapter}{Resumen}
El presente reporte documenta el diseño y desarrollo de un \textbf{metaverso
escolar} para el Tecnológico Nacional de México, entendido como un conjunto de
\emph{prácticas} que el estudiante realiza dentro de un videojuego desarrollado
en Unreal Engine. El núcleo del proyecto es un \textbf{servidor web} construido
con Laravel y PostgreSQL que administra usuarios, materias, grupos, prácticas,
agenda y calificaciones, y expone una \textbf{API} que el juego consume.

Se implementaron tres grandes componentes: (1) el \textbf{backend y la API} con
autenticación mediante \emph{magic link} de un solo uso; (2) un \textbf{panel
para docentes} que permite consultar grupos, generar accesos y revisar
calificaciones; y (3) la \textbf{integración con Moodle mediante el estándar LTI
1.3}, que permite que el alumno ingrese desde una actividad de Moodle (con
identidad automática) y que la calificación obtenida en el juego regrese de
forma automática al libro de calificaciones (servicio AGS).

El sistema se construyó siguiendo prácticas de ingeniería de software modernas:
desarrollo guiado por pruebas (TDD), revisión por pares de cada incremento y una
suite automatizada de \textbf{59 pruebas}. Como resultado, se obtuvo una
plataforma funcional, segura y verificada de extremo a extremo contra una
instancia real de Moodle.

\medskip
\noindent\textbf{Palabras clave:} LTI 1.3, Moodle, Laravel, PostgreSQL, Unreal
Engine, API REST, magic link, AGS, calificaciones, metaverso educativo.

\tableofcontents
\listoffigures
\lstlistoflistings

\clearpage
\pagenumbering{arabic}

% =========================== CAP 1: INTRODUCCIÓN ===========================
\chapter{Introducción}

La transformación digital de la educación ha impulsado el uso de entornos
inmersivos y gamificados como apoyo al aprendizaje. Un \emph{metaverso escolar}
plantea que el estudiante no solo lea o vea contenidos, sino que \textbf{realice
actividades} dentro de un espacio virtual interactivo. Sin embargo, para que una
herramienta así sea adoptable en una institución pública como el Tecnológico
Nacional de México, debe \textbf{integrarse con los sistemas existentes} —
particularmente con la plataforma Moodle que ya utilizan docentes y alumnos— y
hacerlo con la menor fricción administrativa posible.

Este proyecto aborda precisamente ese reto: construir la \textbf{infraestructura
de software} (servidor, base de datos, API y mecanismos de autenticación) que
permita conectar un juego desarrollado en Unreal Engine con la gestión académica
real, incluyendo el \textbf{retorno automático de calificaciones} a Moodle.

\section{Contexto}
El proyecto se desarrolla para el TecNM, institución que emplea Moodle como
sistema de gestión del aprendizaje (LMS). Los docentes administran sus cursos en
Moodle y registran calificaciones en él. Cualquier herramienta nueva que
pretenda ser usada de forma cotidiana debe respetar ese flujo de trabajo.

\section{Descripción del problema}
No existía un mecanismo que permitiera: (a) que un alumno ingresara a un juego
educativo identificándose automáticamente con su identidad institucional; (b)
que el resultado obtenido en el juego se registrara como calificación; ni (c)
que el docente pudiera operar todo esto sin depender de tareas manuales o de
intervención constante del administrador del sistema.

\section{Estructura del documento}
El capítulo 2 presenta la justificación; el 3, los objetivos; el 4 describe el
entorno del proyecto; el 5 establece alcances y limitaciones; el 6 desarrolla el
marco teórico; el 7 detalla el desarrollo técnico (la mayor parte del trabajo);
el 8 expone los resultados; el 9, las conclusiones; y el 10, las competencias
desarrolladas. Al final se incluyen referencias y anexos con código.

% =========================== CAP 2: JUSTIFICACIÓN ===========================
\chapter{Justificación}

La incorporación de un metaverso educativo aporta motivación y aprendizaje
activo, pero su valor real depende de que se integre con la operación académica.
Las razones que justifican este proyecto son:

\begin{itemize}[leftmargin=1.4em]
  \item \textbf{Reducción de carga administrativa:} al devolver las
  calificaciones automáticamente a Moodle, el docente evita capturar notas a
  mano, disminuyendo errores y tiempo invertido.
  \item \textbf{Mínima burocracia de adopción:} el uso del estándar LTI 1.3
  permite que el administrador de Moodle registre la herramienta \emph{una sola
  vez}; después, cualquier docente la usa como una actividad más.
  \item \textbf{Identidad sin contraseñas adicionales:} el alumno no crea una
  cuenta nueva; ingresa desde Moodle o mediante un enlace seguro de un solo uso.
  \item \textbf{Reutilización y escalabilidad:} la API es agnóstica del motor de
  juego; hoy la consume Unreal Engine, pero podría consumirla cualquier cliente.
  \item \textbf{Seguridad:} los tokens se almacenan cifrados (hash), son de un
  solo uso y con expiración, y la comunicación con Moodle se firma
  criptográficamente.
\end{itemize}

% =========================== CAP 3: OBJETIVOS ===========================
\chapter{Objetivos}

\section{Objetivo general}
Diseñar y desarrollar la infraestructura de software (backend, API, panel
docente e integración con Moodle vía LTI 1.3) que permita a los estudiantes del
TecNM realizar prácticas dentro de un videojuego y que sus calificaciones se
registren automáticamente en la plataforma institucional.

\section{Objetivos específicos}
\begin{enumerate}[leftmargin=1.6em]
  \item Modelar e implementar la base de datos académica (usuarios, roles,
  carreras, materias, grupos, prácticas, agenda y sesiones) en PostgreSQL.
  \item Implementar un esquema de autenticación por \emph{magic link} de un solo
  uso para lanzar el juego.
  \item Desarrollar la API REST que el juego consume para canjear el acceso,
  iniciar la sesión de práctica y registrar la calificación.
  \item Construir un panel web para docentes con control de acceso por rol.
  \item Integrar el sistema con Moodle mediante LTI 1.3 (OIDC, Deep Linking y
  AGS), con aprovisionamiento automático del alumno.
  \item Verificar el sistema mediante una suite de pruebas automatizadas y una
  prueba de extremo a extremo contra una instancia real de Moodle.
\end{enumerate}

% =========================== CAP 4: ENTORNO ===========================
\chapter{Caracterización del proyecto}

\section{Área de aplicación}
El proyecto se ubica en el área de \emph{tecnologías educativas} y desarrollo de
software para instituciones de educación superior. Involucra desarrollo backend,
diseño de bases de datos, seguridad, integración de sistemas (LMS) y
comunicación con un motor de videojuegos.

\section{Herramientas y tecnologías}
\begin{longtable}{p{4.2cm} p{9.5cm}}
\toprule
\textbf{Tecnología} & \textbf{Uso en el proyecto} \\
\midrule
\endhead
Laravel 12 (PHP 8.3+) & Framework del backend y la API REST. \\
PostgreSQL & Base de datos relacional (incluye columnas \texttt{jsonb}). \\
Laravel Sanctum & Emisión de tokens \emph{Bearer} para la sesión del juego. \\
\texttt{packbackbooks/lti-1p3-tool} & Librería de integración LTI 1.3 (OIDC, JWT, Deep Linking, AGS). \\
Moodle 5.0 & LMS institucional con el que se integra (vía LTI 1.3). \\
Docker & Ejecución de una instancia local de Moodle para pruebas. \\
Unreal Engine (C++) & Cliente del juego que consume la API. \\
Pest / PHPUnit & Pruebas automatizadas. \\
\bottomrule
\end{longtable}

% =========================== CAP 5: ALCANCES ===========================
\chapter{Alcances y limitaciones}

\section{Alcances}
\begin{itemize}[leftmargin=1.4em]
  \item Base de datos académica completa y migraciones versionadas.
  \item Autenticación por magic link (hash, un solo uso, expiración configurable).
  \item API del juego: canje de acceso, inicio y cierre de sesión con calificación
  y telemetría.
  \item Panel docente: vista de grupos y alumnos, generación de enlaces por grupo
  y consulta de resultados.
  \item Integración LTI 1.3 completa: \emph{login} OIDC, \emph{launch}, Deep
  Linking (selección de práctica) y AGS (retorno de calificación).
  \item Aprovisionamiento automático del alumno a partir del lanzamiento de Moodle.
  \item Cliente de ejemplo en Unreal Engine (C++) y suite de 59 pruebas.
\end{itemize}

\section{Limitaciones}
\begin{itemize}[leftmargin=1.4em]
  \item El proyecto de videojuego en Unreal Engine se entrega como cliente de
  ejemplo; el desarrollo artístico/jugable completo queda fuera del alcance.
  \item La integración se validó contra una instancia local de Moodle en Docker;
  el endurecimiento para producción (HTTPS, secretos gestionados, alta
  disponibilidad) se documenta como trabajo futuro.
  \item Servicios LTI opcionales como NRPS (sincronización de listas) y el manejo
  multi-plataforma (varios Moodle simultáneos) se dejan para una fase posterior.
\end{itemize}

% =========================== CAP 6: MARCO TEÓRICO ===========================
\chapter{Marco teórico}

\section{Sistemas de gestión del aprendizaje y Moodle}
Un \emph{Learning Management System} (LMS) es una plataforma para administrar
cursos, materiales, actividades y calificaciones. Moodle es un LMS de código
abierto ampliamente adoptado en el ámbito educativo, incluido el TecNM. Permite
extender su funcionalidad mediante \emph{actividades externas} que se comunican
con herramientas de terceros a través de estándares de interoperabilidad.

\section{El estándar LTI 1.3}
\emph{Learning Tools Interoperability} (LTI) es el estándar de 1EdTech (antes IMS
Global) para conectar un LMS (\emph{plataforma}) con una herramienta externa
(\emph{tool}). La versión 1.3 se apoya en estándares de seguridad modernos:

\begin{itemize}[leftmargin=1.4em]
  \item \textbf{OpenID Connect (OIDC):} flujo de inicio que previene ataques de
  inyección de \emph{login}. El navegador es redirigido entre plataforma y
  herramienta con un parámetro \texttt{state} y un \texttt{nonce}.
  \item \textbf{JWT firmado (JWS):} el mensaje de lanzamiento (\emph{launch}) es
  un \emph{JSON Web Token} firmado por la plataforma con su llave privada; la
  herramienta lo valida con la llave pública de la plataforma (JWKS).
  \item \textbf{Deep Linking:} permite que el docente, al crear la actividad,
  seleccione \emph{qué} contenido (en este caso, qué práctica) se enlazará.
  \item \textbf{AGS (Assignment and Grade Services):} servicios que permiten a la
  herramienta \textbf{enviar calificaciones} de regreso al libro de
  calificaciones del LMS.
\end{itemize}

La confianza mutua se establece con un \textbf{registro único}: la plataforma y
la herramienta intercambian sus URLs y un \texttt{client\_id}, y cada una publica
sus llaves públicas en un \emph{JSON Web Key Set} (JWKS).

\section{JSON Web Tokens y JWKS}
Un JWT es una cadena con tres partes (cabecera, cuerpo y firma) codificadas en
Base64URL. La firma garantiza integridad y autenticidad. Para validarla se
requiere la llave pública correspondiente, que se publica en un JWKS: un JSON con
las llaves (módulo \texttt{n}, exponente \texttt{e}, identificador \texttt{kid} y
algoritmo \texttt{alg}). En este proyecto la herramienta genera su propio par de
llaves RSA y expone su JWKS para que Moodle valide los mensajes firmados por la
herramienta (respuesta de Deep Linking y peticiones de AGS).

\section{Laravel, PostgreSQL y Sanctum}
Laravel es un framework PHP que favorece el desarrollo expresivo y seguro
(migraciones, ORM Eloquent, validación, middleware). PostgreSQL es un gestor de
base de datos relacional robusto, con soporte de tipos avanzados como
\texttt{jsonb}, usado aquí para almacenar la telemetría cruda del juego. Laravel
Sanctum provee la emisión de tokens \emph{Bearer} para autenticar las llamadas
de la API que realiza el juego.

\section{Autenticación por \emph{magic link}}
Un \emph{magic link} es un enlace que contiene un token y permite el acceso sin
contraseña. Para que sea seguro: el token viaja solo en el enlace, en la base de
datos se almacena únicamente su \textbf{hash} (SHA-256), es de \textbf{un solo
uso} y tiene \textbf{expiración}. Este patrón reduce el soporte por contraseñas
olvidadas y es ideal para lanzar una aplicación de escritorio (el juego) desde un
enlace.

\section{Unreal Engine y \emph{deep links}}
Unreal Engine es un motor de videojuegos. Para que un enlace en el navegador
abra el juego de escritorio se utiliza un \emph{deep link} con un esquema propio
(por ejemplo \texttt{tecnm-metaverso://}). El backend permanece agnóstico del
transporte: solo entrega y valida tokens; el juego, mediante el módulo HTTP de
Unreal, consume la API.

% =========================== CAP 7: DESARROLLO ===========================
\chapter{Desarrollo}

Este capítulo describe la solución construida. Se organiza en: arquitectura
general, modelo de datos, autenticación por magic link, API del juego, panel
docente, integración LTI 1.3, cliente Unreal y metodología de trabajo.

\section{Arquitectura general}
La solución se compone de tres actores principales y un LMS:

\begin{figure}[h]
\centering
\begin{tikzpicture}[
  node distance=0.9cm and 1.6cm,
  caja/.style={rectangle,rounded corners,draw=azultec,thick,fill=azultec!8,
    text width=3.1cm,align=center,minimum height=1.1cm,font=\small},
  flecha/.style={-{Stealth[length=3mm]},thick,azultec}]
  \node[caja] (moodle) {Moodle\\(LMS / plataforma LTI)};
  \node[caja,right=of moodle] (backend) {Backend Laravel\\+ PostgreSQL\\(API)};
  \node[caja,below=of backend] (juego) {Juego Unreal\\(cliente)};
  \node[caja,below=of moodle] (docente) {Panel docente\\(web)};
  \draw[flecha] (moodle) -- node[above,font=\scriptsize]{LTI 1.3} (backend);
  \draw[flecha] (backend) -- node[right,font=\scriptsize]{API REST} (juego);
  \draw[flecha] (docente) -- node[left,font=\scriptsize]{sesión web} (backend);
  \draw[flecha,dashed] (backend) to[bend right=18] node[below,font=\scriptsize]{AGS (calificación)} (moodle);
\end{tikzpicture}
\caption{Arquitectura general del sistema.}
\label{fig:arq}
\end{figure}

El backend es el \emph{cerebro}: ninguna persona lo usa directamente; lo
consultan el juego (vía API) y el panel (vía web), y se comunica con Moodle por
LTI. Esta separación de responsabilidades permite evolucionar cada parte de forma
independiente.

\section{Modelo de datos}
El esquema se diseñó normalizado, con nombres de tabla en español y llaves
primarias propias (por ejemplo \texttt{id\_usuario}). La
Tabla~\ref{tab:tablas} resume las entidades principales.

\begin{longtable}{p{3.6cm} p{10cm}}
\caption{Entidades principales del modelo de datos.}\label{tab:tablas}\\
\toprule
\textbf{Tabla} & \textbf{Propósito} \\
\midrule
\endfirsthead
\toprule \textbf{Tabla} & \textbf{Propósito} \\ \midrule
\endhead
\texttt{roles}, \texttt{usuarios} & Identidad y rol (Alumno, Maestro, Coordinador, Admin). \\
\texttt{alumnos}, \texttt{maestros} & Datos específicos de cada tipo de usuario. \\
\texttt{carreras}, \texttt{materias}, \texttt{materia\_carrera} & Catálogo académico y su relación. \\
\texttt{ciclos\_escolares}, \texttt{grupos}, \texttt{inscripciones} & Organización por ciclo y grupo. \\
\texttt{practicas} & Actividad a realizar en el juego (incluye escena de referencia). \\
\texttt{eventos\_agenda} & Programación de una práctica para un grupo. \\
\texttt{sesiones\_practica} & Sesión jugada: estatus, calificación y telemetría (\texttt{jsonb}). \\
\texttt{tokens\_juego} & Magic link: \emph{hash} del token, expiración, uso. \\
\bottomrule
\end{longtable}

Se aplicaron correcciones de seguridad respecto al diseño inicial: el token del
magic link se guarda \textbf{hasheado} (no en texto plano), y la sesión de
práctica admite que la columna \texttt{id\_evento} sea nula para los lanzamientos
por LTI (que se anclan directamente a una práctica).

\section{Autenticación por magic link}
El servicio de magic link genera un token aleatorio robusto, almacena su hash
SHA-256 y construye tanto la URL web como el \emph{deep link} para el juego. La
vigencia es configurable (120 minutos por omisión). El
Listado~\ref{lst:magic} muestra el núcleo del servicio.

\begin{lstlisting}[language=PHP,caption={Generación del magic link (extracto).},label=lst:magic]
public function generar(int $idUsuario, int $idEvento, string $plataforma = 'unreal'): array {
    $token = Str::random(64);
    $ttl   = (int) config('metaverso.magic_link_ttl_minutes', 120);

    $modelo = TokenJuego::create([
        'id_usuario'       => $idUsuario,
        'id_evento'        => $idEvento,
        'token_hash'       => $this->hash($token),   // solo se guarda el hash
        'plataforma'       => $plataforma,
        'fecha_expiracion' => now()->addMinutes($ttl),
        'usado'            => false,
    ]);

    $scheme = config('metaverso.deeplink_scheme', 'tecnm-metaverso');
    return [
        'token'    => $token,
        'modelo'   => $modelo,
        'url'      => URL::to('/jugar/'.$token),
        'deeplink' => "{$scheme}://play?token={$token}",
    ];
}
\end{lstlisting}

\section{API del juego}
El juego consume una API bajo el prefijo \texttt{/api/game}. El canje del token
se valida de forma \textbf{atómica} para impedir su reutilización en condiciones
de concurrencia. La Tabla~\ref{tab:api} resume los puntos de acceso.

\begin{longtable}{p{1.3cm} p{5.4cm} p{6.4cm}}
\caption{Puntos de acceso de la API del juego.}\label{tab:api}\\
\toprule
\textbf{Verbo} & \textbf{Ruta} & \textbf{Propósito} \\ \midrule
\endhead
POST & \texttt{/api/game/redeem} & Canjea el magic link y entrega un Bearer. \\
GET & \texttt{/api/game/me} & Datos del alumno autenticado. \\
POST & \texttt{/api/game/sessions} & Inicia una sesión de práctica. \\
POST & \texttt{/api/game/sessions/\{id\}/complete} & Registra calificación (0--100) y telemetría. \\
POST & \texttt{/api/game/lti-redeem} & Canje específico para sesiones creadas por LTI. \\
\bottomrule
\end{longtable}

Las respuestas de error son consistentes (401 token inválido, 410 expirado o
usado, 422 validación, 403 sesión ajena, 409 estado incorrecto). Al completar una
sesión se almacena la calificación y la telemetría en \texttt{jsonb}; si la
sesión proviene de LTI, además se dispara el envío de la nota a Moodle.

\section{Panel docente}
El panel permite al docente operar sin consola ni base de datos. El acceso se
realiza mediante una \textbf{URL firmada temporal} de Laravel (sin contraseña ni
servidor de correo), generada con un comando de consola. El inicio de sesión del
panel se encapsuló en un \textbf{único método} reutilizable, de modo que la futura
entrada vía LTI lo aproveche sin cambios. Las funciones principales son:

\begin{itemize}[leftmargin=1.4em]
  \item \textbf{Tablero:} grupos del docente (o todos, si es Coordinador/Admin).
  \item \textbf{Detalle de grupo:} alumnos inscritos y eventos/prácticas.
  \item \textbf{Generación de enlaces:} un magic link por alumno del grupo, con
  exportación a CSV (generada en el cliente para no regenerar tokens).
  \item \textbf{Resultados:} sesiones por práctica con calificación y telemetría.
\end{itemize}

La autorización es estricta: un docente solo accede a sus propios grupos, y el
rol Alumno no puede ingresar al panel.

\section{Integración LTI 1.3 con Moodle}
Es el componente más extenso. Permite que el alumno entre desde Moodle y que la
calificación regrese automáticamente. El flujo completo se ilustra en la
Figura~\ref{fig:lti}.

\begin{figure}[h]
\centering
\begin{tikzpicture}[
  paso/.style={rectangle,rounded corners,draw=azultec,thick,fill=azultec!6,
    text width=11.5cm,align=left,font=\small,inner sep=6pt},
  node distance=0.45cm]
  \node[paso] (p1) {\textbf{1.} El docente agrega la actividad y, vía \emph{Deep Linking}, elige la práctica.};
  \node[paso,below=of p1] (p2) {\textbf{2.} El alumno hace clic en la actividad; Moodle inicia el \emph{login} OIDC.};
  \node[paso,below=of p2] (p3) {\textbf{3.} Moodle envía el \emph{launch} (JWT firmado) al backend, que lo valida.};
  \node[paso,below=of p3] (p4) {\textbf{4.} El backend aprovisiona al alumno y crea la sesión de práctica.};
  \node[paso,below=of p4] (p5) {\textbf{5.} Se muestra ``Abrir juego''; Unreal canjea el token y obtiene un Bearer.};
  \node[paso,below=of p5] (p6) {\textbf{6.} Al terminar, el juego envía la calificación; el backend la reenvía a Moodle (AGS).};
  \draw[-{Stealth},azultec,thick] (p1)--(p2);
  \draw[-{Stealth},azultec,thick] (p2)--(p3);
  \draw[-{Stealth},azultec,thick] (p3)--(p4);
  \draw[-{Stealth},azultec,thick] (p4)--(p5);
  \draw[-{Stealth},azultec,thick] (p5)--(p6);
\end{tikzpicture}
\caption{Flujo de integración LTI 1.3 (Moodle $\rightarrow$ juego $\rightarrow$ calificación).}
\label{fig:lti}
\end{figure}

\subsection{Puntos de acceso LTI}
\begin{itemize}[leftmargin=1.4em]
  \item \texttt{GET /lti/jwks}: publica la llave pública de la herramienta.
  \item \texttt{GET|POST /lti/login}: inicio OIDC; redirige a la plataforma.
  \item \texttt{POST /lti/launch}: valida el lanzamiento, aprovisiona al alumno y
  crea la sesión, o (si es Deep Linking) muestra el selector de práctica.
  \item \texttt{GET|POST /lti/deeplink}: selector de práctica y respuesta firmada.
\end{itemize}

\subsection{Aprovisionamiento automático}
Cuando un alumno entra por primera vez, el backend lo identifica por su
\texttt{lti\_user\_id} (proporcionado por Moodle) y, si no existe, crea su usuario
y registro de alumno de forma \textbf{idempotente}. Así se elimina el alta manual
de estudiantes.

\subsection{Retorno de calificación (AGS)}
Al completar la sesión, si esta proviene de LTI, el backend envía la calificación
al \emph{line item} de Moodle mediante AGS, firmando la petición con su llave
privada. El envío se realiza de forma resiliente: si Moodle no responde, el error
se registra sin interrumpir el flujo del alumno.

\subsection{Consideraciones de despliegue local}
Para las pruebas se levantó Moodle 5.0 en Docker. Se resolvieron los aspectos
clásicos de un entorno LTI local sobre HTTP: el uso de \emph{cookies same-site}
para el parámetro \texttt{state}, la habilitación del puerto del backend en la
seguridad de cURL de Moodle, y una \textbf{tolerancia de reloj} (\emph{leeway})
en la validación de los JWT para absorber el desfase de reloj del contenedor.

\section{Cliente en Unreal Engine}
Se desarrolló un cliente de ejemplo en C++ que, usando el módulo HTTP nativo de
Unreal, canjea el token y envía la calificación. Para integrarlo a la jugabilidad
se proveyó además un \emph{GameInstanceSubsystem} con funciones invocables desde
Blueprint: \texttt{IniciarSesion}, \texttt{SumarAcierto} y \texttt{Terminar}. El
Listado~\ref{lst:unreal} muestra el envío de la calificación.

\begin{lstlisting}[language=C++,caption={Envío de la calificación desde Unreal (extracto).},label=lst:unreal]
void UMetaversoSubsystem::EnviarCalificacion(float Calificacion) {
    TSharedRef<IHttpRequest, ESPMode::ThreadSafe> Req = FHttpModule::Get().CreateRequest();
    Req->SetURL(FString::Printf(TEXT("%s/api/game/sessions/%d/complete"), *BaseUrl, IdSesion));
    Req->SetVerb(TEXT("POST"));
    Req->SetHeader(TEXT("Content-Type"), TEXT("application/json"));
    Req->SetHeader(TEXT("Authorization"), TEXT("Bearer ") + AccessToken);
    Req->SetContentAsString(FString::Printf(
        TEXT("{\"calificacion\":%.1f,\"datos_resultado\":{\"fuente\":\"unreal\"}}"), Calificacion));
    Req->ProcessRequest();
}
\end{lstlisting}

\section{Metodología de desarrollo}
El desarrollo siguió un proceso disciplinado: para cada incremento se escribió
primero la prueba (TDD), luego la implementación mínima, y finalmente una
\textbf{revisión} de cumplimiento y calidad antes de integrarse. El trabajo se
dividió en sub-proyectos (backend, panel, LTI), cada uno con su especificación,
plan de implementación y suite de pruebas, integrándose mediante control de
versiones (Git) con revisiones por incremento.

% =========================== CAP 8: RESULTADOS ===========================
\chapter{Resultados}

Se obtuvo una plataforma funcional y verificada. Los principales resultados son:

\begin{itemize}[leftmargin=1.4em]
  \item \textbf{Backend y API} operativos, con autenticación por magic link
  segura (hash, un solo uso, expiración) y registro de calificación y telemetría.
  \item \textbf{Panel docente} con acceso por enlace firmado, control por rol,
  generación de enlaces por grupo (con CSV) y consulta de resultados.
  \item \textbf{Integración LTI 1.3} completa, probada de extremo a extremo contra
  Moodle 5.0: el alumno entra desde la actividad, juega y la \textbf{calificación
  regresa automáticamente} al libro de calificaciones.
  \item \textbf{Suite de pruebas} con 59 casos automatizados en verde, que cubren
  el flujo de tokens, sesiones, autorización del panel y la lógica de LTI.
  \item \textbf{Documentación} técnica y operativa: guías de configuración de
  Moodle, de uso del panel y de integración del cliente Unreal.
\end{itemize}

La Tabla~\ref{tab:result} resume el estado de los objetivos específicos.

\begin{longtable}{p{10.5cm} p{2.8cm}}
\caption{Cumplimiento de objetivos específicos.}\label{tab:result}\\
\toprule
\textbf{Objetivo} & \textbf{Estado} \\ \midrule
\endhead
Modelo de datos académico en PostgreSQL & Cumplido \\
Autenticación por magic link de un solo uso & Cumplido \\
API del juego (canje, sesión, calificación) & Cumplido \\
Panel docente con control por rol & Cumplido \\
Integración LTI 1.3 (OIDC, Deep Linking, AGS) & Cumplido \\
Pruebas automatizadas y prueba de extremo a extremo & Cumplido \\
\bottomrule
\end{longtable}

% =========================== CAP 9: CONCLUSIONES ===========================
\chapter{Conclusiones y recomendaciones}

\section{Conclusiones}
El proyecto demuestra que es posible integrar un metaverso educativo con la
operación académica real de una institución pública, con baja fricción de
adopción y respetando el flujo de trabajo de los docentes en Moodle. El uso del
estándar LTI 1.3 fue determinante: con un único registro de la herramienta, el
sistema queda disponible para todos los docentes, y la calificación regresa
automáticamente, eliminando captura manual.

Desde el punto de vista de ingeniería, el desarrollo guiado por pruebas y la
revisión por incrementos resultaron en un sistema confiable (59 pruebas) y
mantenible, con responsabilidades bien separadas.

\section{Recomendaciones y trabajo futuro}
\begin{itemize}[leftmargin=1.4em]
  \item \textbf{Producción:} servir el backend sobre HTTPS con certificados
  válidos, gestionar secretos de forma segura y proteger el endpoint interno de
  generación de enlaces con autenticación de docente.
  \item \textbf{Panel ampliado:} permitir crear prácticas y agendar eventos desde
  el panel, y mostrar las sesiones LTI en la vista de resultados.
  \item \textbf{LTI avanzado:} incorporar NRPS (sincronización de listas) y
  soporte multi-plataforma.
  \item \textbf{Juego:} desarrollar el contenido jugable en Unreal y registrar el
  esquema de \emph{deep link} para abrir el juego automáticamente desde el enlace.
\end{itemize}

% =========================== CAP 10: COMPETENCIAS ===========================
\chapter{Competencias desarrolladas}
\begin{itemize}[leftmargin=1.4em]
  \item Diseño e implementación de bases de datos relacionales normalizadas.
  \item Desarrollo de APIs REST seguras y su autenticación por tokens.
  \item Integración de sistemas mediante estándares (LTI 1.3, OIDC, JWT, OAuth2).
  \item Aplicación de criptografía práctica (hashing, firma de JWT, JWKS).
  \item Desarrollo guiado por pruebas y aseguramiento de la calidad.
  \item Uso de contenedores (Docker) y diagnóstico de problemas de integración.
  \item Comunicación con un motor de videojuegos (Unreal Engine, C++ y HTTP).
  \item Documentación técnica profesional.
\end{itemize}

% =========================== REFERENCIAS ===========================
\chapter*{Referencias}
\addcontentsline{toc}{chapter}{Referencias}
\begin{enumerate}[leftmargin=1.6em]
  \item 1EdTech Consortium. \emph{Learning Tools Interoperability (LTI) Core
  Specification v1.3}. Disponible en: \url{https://www.imsglobal.org/spec/lti/v1p3}
  \item 1EdTech Consortium. \emph{LTI Assignment and Grade Services (AGS) v2.0}.
  Disponible en: \url{https://www.imsglobal.org/spec/lti-ags/v2p0}
  \item Moodle. \emph{Documentación oficial}. Disponible en:
  \url{https://docs.moodle.org/}
  \item Laravel. \emph{Documentation}. Disponible en: \url{https://laravel.com/docs}
  \item Laravel Sanctum. Disponible en: \url{https://laravel.com/docs/sanctum}
  \item PostgreSQL Global Development Group. \emph{PostgreSQL Documentation}.
  Disponible en: \url{https://www.postgresql.org/docs/}
  \item Epic Games. \emph{Unreal Engine Documentation}. Disponible en:
  \url{https://dev.epicgames.com/documentation/}
  \item Packback. \emph{lti-1p3-tool (PHP LTI 1.3 library)}. Disponible en:
  \url{https://github.com/packbackbooks/lti-1-3-php-library}
\end{enumerate}

% =========================== ANEXOS ===========================
\appendix
\chapter{Anexos}

\section{Comandos de consola del sistema}
\begin{lstlisting}[language=bash,caption={Comandos artisan del proyecto.}]
# Generar un magic link de prueba (juego):
php artisan metaverso:magic-link <id_usuario> <id_evento>

# Generar un enlace de acceso al panel docente:
php artisan metaverso:panel-acceso <id_usuario>

# Generar el par de llaves RSA de la herramienta LTI:
php artisan metaverso:lti-generar-llaves

# Registrar la plataforma Moodle (LTI):
php artisan metaverso:lti-registrar-plataforma \
  --issuer=... --client-id=... --deployment-id=... \
  --auth-login-url=... --auth-token-url=... --jwks-url=...
\end{lstlisting}

\section{Flujo de la API consumida por el juego}
\begin{lstlisting}[language=bash,caption={Secuencia de llamadas (resumen).}]
1) POST /api/game/redeem        { token }            -> { access_token, alumno, practica }
2) POST /api/game/sessions      { id_evento }         (Bearer)  -> { id_sesion }
3) POST /api/game/sessions/{id}/complete
        { calificacion, datos_resultado }            (Bearer)  -> 200
\end{lstlisting}

\end{document}
